\section{Conclusiones}
\label{sec:conclusiones}

Este trabajo describió la versión de ocho agentes de EAM-TMS, un sistema de memoria
transactiva multi-agente construido sobre la Memoria Asociativa Entrópica ponderada
(W-EAM) y su extensión hetero-asociativa
(EHAM)~\cite{pineda2021entropic,pineda2022weighted,morales2024hetero,morales2025missing}.
El sistema cubre las ocho categorías de ETH-80 y opera de forma bidireccional
texto $\leftrightarrow$ imagen con rechazo claro en ambas direcciones, sin módulos
externos de clasificación ni umbrales supervisados.

\paragraph{Logros principales.}
Sobre el banco de 411 consultas ($N = 400$, 5 semillas), el sistema alcanzó una
\textbf{precisión temprana del 88.0\,\%} y una \textbf{precisión madura del 93.2\,\%}
con la lectura B1 del directorio, con una fidelidad temprana-madura del 93.5\,\% y la
participación del agente dominante en \textbf{13.0\,\%} (ideal 12.5\,\%).
En la dirección inversa, el hemisferio visual enrutó el \textbf{75.0\,\%} de las 656
imágenes de prueba con \textbf{cero falsos enrutamientos} (100\,\% de precisión sobre
las aceptadas) y una evocación de dominio del \textbf{85.3\,\%}; su directorio alcanzó
el balance perfecto: \textbf{3.000 de 3.000 bits} de entropía.
La especificidad del contenido fue \textbf{perfecta} en todo el rango de llenado
probado (0\,\% de falsos aceptados con siete clases contrarias por agente).

\paragraph{Respuesta a la pregunta de investigación.}
La maquinaria transactiva \emph{escala}: directorios, lectura B1, rechazo por
containment, actualización de directorio y coordinación punto a punto pasan de 3 a 8
agentes sin cambios de mecanismo.
Los límites encontrados estaban en las \emph{representaciones}: la binarización por
signo (resuelta con la cuantización por magnitud, que llevó la precisión temprana del
50\,\% al 88\,\%), la separabilidad del espacio latente del encoder (cuyo
reentrenamiento resolvió la dominancia de \texttt{apple}: 20\,\% $\to$ 13.0\,\%), y la
asimetría léxica de la fuente de conocimiento (que fija el techo de las clases con
vocabulario compartido).
El costo de la escala aparece donde la teoría de Wegner lo predice: un grupo con más
especialidades necesita más interacciones, y más distintivas, para formar su directorio.

\paragraph{Lecciones de esta versión.}
Primero, la representación pesa tanto como el mecanismo: dos de los tres límites
observados se resolvieron cambiando representaciones sin tocar las memorias.
Segundo, los sesgos observados en un hemisferio pueden originarse en el otro: la
dominancia ``léxica'' de \texttt{apple} era en buena parte solapamiento latente del
encoder.
Tercero, la fidelidad del pipeline es un resultado en sí mismo: la auditoría que
eliminó el único vector sintético del sistema también corrigió la escala global de
cuantización, e hizo trazables todos los números reportados.

\paragraph{Trabajo futuro.}
Los siguientes pasos más relevantes son:
\begin{enumerate}
    \item \textbf{Doble compuerta de contenido en la fase madura textual}: cerrar la
    fuga de consultas fuera de dominio verificando el containment en el agente destino,
    no solo la señal del directorio.
    \item \textbf{Augmentación más rica del llenado visual}: reducir el 25\,\% de
    rechazo residual de test, que proviene del containment conjunto y no responde a
    $\xi$.
    \item \textbf{Fuentes de conocimiento complementarias a ConceptNet}: elevar el techo
    de las clases con léxico compartido con conocimiento alineado al lenguaje de las
    consultas, único tipo de enriquecimiento que demostró ayudar.
    \item \textbf{Partición leave-object-out}: validar la generalización visual sobre
    objetos físicos completamente no vistos.
    \item \textbf{Aprendizaje en fase madura}: extender el sistema para que los agentes
    sigan almacenando labels, imágenes y asociaciones después de formar el directorio,
    cubriendo \emph{transactive encoding} y \emph{transactive retrieval} en sentido
    amplio.
    \item \textbf{Rotación de miembros}: explorar olvido controlado y reorganización del
    directorio cuando un agente abandona o se incorpora al grupo.
\end{enumerate}

\bigskip

En resumen, los resultados muestran que las Memorias Asociativas Entrópicas de Pineda y
Morales son una base suficiente para construir un sistema de memoria transactiva
multi-agente que escala a ocho especialistas, bidireccional y con rechazo controlado,
sin recurrir a modelos de aprendizaje profundo supervisado para la toma de decisiones;
y que, cuando el escalado encuentra límites, estos viven en las representaciones que
alimentan a las memorias y no en el mecanismo asociativo ni en la maquinaria transactiva.
